\documentclass[../../../include/open-logic-section]{subfiles}

\begin{document}

\olfileid{cmp}{rec}{ore}
\olsection{其他递归方式}

% Part: computability
% Chapter: recursive-functions
% Section: other-recursions

利用配对和序列化，我们可以论证更奇特（且有用）的原始递归形式的合理性。
例如，联立定义两个函数通常很有用，如下面的定义所示：
\begin{align*}
h_0(\vec x, 0) & = f_0(\vec x) \\
h_1(\vec x, 0) & = f_1(\vec x) \\
h_0(\vec x, y+1) & = g_0(\vec x, y, h_0(\vec x, y), h_1(\vec x, y)) \\
h_1(\vec x, y+1) & = g_1(\vec x, y, h_0(\vec x, y), h_1(\vec x, y))
\end{align*}
这就是\emph{联立递归}的一个实例。
定义函数的另一种有用方法是，根据\emph{所有}值 $h(\vec x, 0)$, ……,~$h(\vec x, y)$ 来给出 $h(\vec x, y+1)$ 的值，如下面的定义所示：
\begin{align*}
h(\vec x, 0) & = f(\vec x) \\
h(\vec x, y+1) & = g(\vec x, y, \tuple{h(\vec x, 0), \dots, h(\vec x, y)}).
\end{align*}
下面的模式更简洁地刻画了这一思想：
\[
h(\vec x, y) = g(\vec x, y, \tuple{h(\vec x, 0), \dots, h(\vec x, y-1)})
\]
其中约定当 $y$ 为 $0$ 时，$g$ 的最后一个自变元就是空序列。
无论哪种表述，其思想都是在计算“后继步”时，函数 $h$ 可以利用已计算出的整个值序列。
这被称为\emph{串值递归}。
对于一个具体的例子，它可以用来论证以下类型定义的合理性：
\begin{align*}
h(\vec x, y) & = \begin{cases}
  g(\vec x, y, h(\vec x, k(\vec x, y))) & \text{若 $k(\vec x, y) < y$} \\
  f(\vec x) & \text{否则}
\end{cases}
\end{align*}
换句话说，$h$ 在 $y$ 处的值可以根据 $h$ 在由~$k$ 给出的\emph{任何}先前值上的取值来计算。

\begin{prob}
  用串值递归定义余数函数 $r(x,y)$。（若 $x$, $y$ 是自然数且 $y > 0$，则 $r(x,y)$ 是小于~$y$ 且满足对某个~$z$ 有 $x = z\times y + r(x,y)$ 的数。为明确起见，我们约定若 $y=0$，则 $r(x,0) = 0$。）
\end{prob}

你应该思考如何用普通原始递归得到这些函数。
最后再介绍一种原始递归的版本，它更加灵活，允许在递归过程中改变\emph{参数}（辅助值）：
\begin{align*}
h(\vec x, 0) & = f(\vec x) \\
h(\vec x, y+1) & = g(\vec x, y, h(k(\vec x), y))
\end{align*}
这也能够用普通原始递归来模拟。（这做起来有些技巧。提示一下：尝试手动展开计算过程。）

\end{document}